\documentclass[]{ConspectBook}
\let\enquote\emph

%  1) Charges
%  2) ChargeSystem
%  3) EBMotion
%  4) ElEnergy
%  5) ElMagInduction
%  6) EMOscillations
%  7) EMWaves
%  8) FileldInDielectrics
%  9) FileldInMetall
% 10) MagEnergy
% 11) MagFieldMedia
% 12) MagFieldVacuum
% 13) MaxwellEqns
% 14) Radiation
% 15) Relativity
% 16) SteadyCurrent
% 17) SteadyEfield
% 18) Telegraph
% 19) Waveguides

\newcommand{\includechapter}[1]{\input{#1/#1}}
\newcommand{\includetikz}[2]{\input{#1/tikz/#2.tikz}}


\begin{document}


%=========================================================
\chapter{Магнітні кола}\label{\currfilebase}
\graphicspath{{\currfilebase/Pictures}}
%=========================================================

Досі ми розглядали електричні кола, в яких струм тече по провідниках, а магнітне поле відігравало роль лише супутнього явища --- воно виникало навколо провідників зі струмом, створювало індуктивність, взаємоіндукцію, і через нього працював трансформатор. Проте є цілий клас пристроїв, де магнітне поле стає \emph{основним} робочим середовищем, а його поведінку зручно описувати мовою, що структурно повторює закони електричних кіл. Йдеться про \emphx{магнітні кола} --- системи феромагнітних осердь, котушок і повітряних зазорів, у яких магнітний потік замикається певним чином, подібно до того, як електричний струм замикається в електричному колі.

Ідея цього розділу така: ми побачимо, що між електричним і магнітним колом існує глибока формальна аналогія. Якщо замінити електричний струм на магнітний потік, ЕРС --- на намагнічувальну силу, а опір --- на магнітний опір, то закони, якими ми користувалися для електричних кіл (закон Ома, правила Кірхгофа, послідовне й паралельне з'єднання), збережуть свою форму. Це дозволить розраховувати магнітні системи --- від дроселів і трансформаторів до електричних машин --- тими самими методами, що й звичайні електричні кола.

Але є й принципова відмінність, яку треба тримати в голові з самого початку: магнітний опір феромагнітного матеріалу \emph{нелінійний} --- він залежить від величини потоку, тоді як електричний опір у звичайних умовах сталий. Саме тому магнітні кола розраховують або в лінійному наближенні (для малих відхилень від робочої точки), або графічно --- за кривою намагнічування. Ми розглянемо обидва підходи, а також з'ясуємо, як ці поняття застосовуються в реальних пристроях.

%% --------------------------------------------------------
\section{Магнітний опір і закон Ома для магнітного кола}
%% --------------------------------------------------------

Почнемо з найпростішого випадку --- замкненого феромагнітного осердя (тороїда або броньового сердечника), на якому рівномірно намотано котушку з $N$ витків, по якій тече струм $I$. Магнітний потік $\Phi$ замикається всередині осердя. Для такого кола теорема про циркуляцію магнітного поля (закон повного струму) дає:
\begin{equation}\label{eq:Ampere_law_magnetic}
	\oint\limits_{\Gamma} \mathbf{H} \cdot d\mathbf{l} = \frac{4\pi}{c} N I.
\end{equation}

Права частина --- це \emphx{намагнічувальна сила} (НС), або \emphx{магніторушійна сила} (МРС):
\begin{equation}\label{eq:MMF_definition}
	\mathcal{F} = N I.
\end{equation}

Вона відіграє в магнітному колі ту саму роль, що ЕРС в електричному: це «двигун», який змушує магнітний потік циркулювати по колу. Якщо контур $\Gamma$ вибрати вздовж середньої лінії осердя довжиною $l$, а поле $H$ вважати однорідним, то $\oint H \, dl = H l$, і закон повного струму набирає вигляду:
\begin{equation}\label{eq:HL_NI}
	H l = \frac{4\pi}{c} N I.
\end{equation}

Зв'яжемо тепер $H$ з магнітним потоком. За означенням індукція $B = \Phi/S$, де $S$ --- площа поперечного перерізу осердя. Для феромагнетика $B = \mu H$ (у СГС $\mu$ --- безрозмірна магнітна проникність), тому:
\begin{equation}\label{eq:H_from_Phi}
	H = \frac{B}{\mu} = \frac{\Phi}{\mu S}.
\end{equation}

Підставляючи \eqref{eq:H_from_Phi} в \eqref{eq:HL_NI}, отримуємо:
\begin{equation}\label{eq:magnetic_Ohm_derivation}
	\frac{\Phi}{\mu S} \, l = \frac{4\pi}{c} N I
	\quad\Rightarrow\quad
	\Phi = \frac{4\pi N I / c}{\dfrac{l}{\mu S}}.
\end{equation}

Величина в знаменнику
\begin{equation}\label{eq:magnetic_reluctance}
	\tcbhighmath{
		R_{\text{м}} = \frac{l}{\mu S}
	}
\end{equation}
називається \emphx{магнітним опором} (або \emphx{релюктансом}) ділянки магнітного кола. Порівняйте цю формулу з електричним опором $R = \rho l / S = l/(\sigma S)$: структура та сама --- довжина, поділена на «провідність» матеріалу й площу перерізу. Роль питомої провідності $\sigma$ відіграє величина $\mu$ --- \emphx{магнітна проникність}.

Отже, ми отримали \emphx{закон Ома для магнітного кола}:
\begin{equation}\label{eq:Ohm_magnetic}
	\tcbhighmath{
		\Phi = \frac{4\pi \mathcal{F}}{c\,R_{\text{м}}},
	}
\end{equation}
де $\mathcal{F} = N I$ --- намагнічувальна сила, $R_{\text{м}}$ --- магнітний опір, $\Phi$ --- магнітний потік. Формально це той самий закон Ома $I = \emf/R$, лише з іншими позначеннями й множником $4\pi/c$, що виникає через вибір системи одиниць (СГС). Це не випадковий збіг: обидва закони випливають із теореми про циркуляцію відповідного поля (електричного --- для ЕРС, магнітного --- для НС) і означення потокової величини (струму --- як потоку заряду, магнітного потоку --- як потоку індукції).

\begin{Attention}
	Зверніть увагу на важливу відмінність між електричним і магнітним опором. Електричний опір мідного дроту --- величина стала: вона не залежить від того, який струм по ньому тече. Магнітний опір феромагнітного осердя --- \emph{нелінійний}: він залежить від величини потоку $\Phi$, бо проникність $\mu$ залежить від напруженості поля $H$. При малих полях $\mu$ мала, при середніх --- різко зростає, а при великих --- насичується, і $\mu$ падає. Саме тому закон Ома для магнітного кола \eqref{eq:Ohm_magnetic} у строгому сенсі справедливий лише для лінійної ділянки кривої намагнічування або для малих відхилень від робочої точки. Для точних розрахунків доводиться користуватися графічними методами або чисельним моделюванням.
\end{Attention}

%% --------------------------------------------------------
\section{Правила Кірхгофа для магнітного кола}
%% --------------------------------------------------------

Оскільки магнітне коло описується тими самими співвідношеннями, що й електричне, природно очікувати, що й правила Кірхгофа збережуть свою форму. Так і є. Сформулюємо їх.

\paragraph*{Перше правило (для вузлів).} У будь-якій точці розгалуження магнітного кола сума магнітних потоків, що входять у вузол, дорівнює сумі потоків, що виходять:
\begin{equation}\label{eq:Kirchhoff_magnetic_current}
	\tcbhighmath{
		\sum_k \Phi_k = 0.
	}
\end{equation}
Це пряме відображення того факту, що лінії магнітної індукції неперервні: магнітний потік не може «зникати» або «народжуватися» у вузлі --- він лише перерозподіляється між гілками. Аналогія з першим правилом Кірхгофа для струмів повна.

\paragraph*{Друге правило (для контурів).} Для будь-якого замкненого контуру в магнітному колі алгебраїчна сума спадів магнітної напруги на ділянках дорівнює алгебраїчній сумі намагнічувальних сил:
\begin{equation}\label{eq:Kirchhoff_magnetic_voltage}
	\tcbhighmath{
		\sum_k \Phi_k R_{\text{м}k} = \frac{4\pi}{c}\sum_i \mathcal{F}_i.
	}
\end{equation}
Тут $\Phi_k$ --- потік через $k$-ту ділянку з магнітним опором $R_{\text{м}k}$, а $\mathcal{F}_i = N_i I_i$ --- намагнічувальні сили котушок, що входять у вибраний контур. Знаки вибираються за тим самим правилом, що й для електричних кіл: напрямок обходу контуру задає додатний напрямок потоків і НС.

Порівняйте \eqref{eq:Kirchhoff_magnetic_voltage} з другим правилом Кірхгофа для змінного струму: $\sum I_k Z_k = \sum \emf_i$. Форма та сама, лише замість імпедансів стоять магнітні опори, а замість ЕРС --- намагнічувальні сили (з тим самим множником $4\pi/c$, що й у законі Ома для магнітного кола).

\begin{Example}{Розрахунок розгалуженого магнітного кола}
	Розглянемо броньовий сердечник (рис.~\ref{tikz:magnetic_circuit}), у якого центральний стрижень має переріз $S_1$, а два крайніх --- $S_2 = S_1/2$ кожен. На центральному стрижні намотано котушку з $N$ витків, по якій тече струм $I$. Знайдемо магнітний потік у кожній гілці, вважаючи осердя лінійним (проникність $\mu$ стала).

	\begin{center}
		\begin{circuitikz}[scale=1.2, thick]
			% Броньовий сердечник (спрощено)
			\draw[fill=gray!20] (-2,-1.5) rectangle (2,1.5);
			\draw[fill=white] (-1.2,-0.8) rectangle (1.2,0.8);
			\draw[fill=gray!20] (-0.4,-0.8) rectangle (0.4,0.8);
			% Котушка
			\draw[red, thick] (-0.6,0.9) -- (-0.6,-0.9);
			\draw[red, thick] (0.6,0.9) -- (0.6,-0.9);
			\node[red] at (0, 1.1) {$N$ витків};
			\node at (0, 0) {$I$};
			% Позначення потоків
			\draw[->, blue, thick] (0, 0.4) -- (0, 0.9) node[right] {$\Phi_1$};
			\draw[->, blue, thick] (-1.6, 0) -- (-1.6, 0.5) node[left] {$\Phi_2$};
			\draw[->, blue, thick] (1.6, 0) -- (1.6, 0.5) node[right] {$\Phi_3$};
		\end{circuitikz}
		\captionof{figure}{Розгалужене магнітне коло: броньовий сердечник із котушкою на центральному стрижні.}
		\label{tikz:magnetic_circuit}
	\end{center}

	Позначимо довжину середньої лінії центрального стрижня $l_1$, а крайніх --- $l_2$. Магнітні опори:
	\[
		R_{\text{м}1} = \frac{l_1}{\mu S_1}, \qquad
		R_{\text{м}2} = \frac{l_2}{\mu S_2} = \frac{2 l_2}{\mu S_1}.
	\]
	За першим правилом Кірхгофа для вузла, де потік $\Phi_1$ розгалужується на $\Phi_2$ і $\Phi_3$:
	\[
		\Phi_1 = \Phi_2 + \Phi_3.
	\]
	Через симетрію $\Phi_2 = \Phi_3 = \Phi_1/2$. За другим правилом Кірхгофа для контуру, що проходить через центральний стрижень і одне з бічних відгалужень:
	\[
		\Phi_1 R_{\text{м}1} + \Phi_2 R_{\text{м}2} = \frac{4\pi}{c} N I.
	\]
	Підставляючи $\Phi_2 = \Phi_1/2$, отримуємо:
	\[
		\Phi_1 \left(R_{\text{м}1} + \frac{R_{\text{м}2}}{2}\right) = \frac{4\pi}{c} N I
		\quad\Rightarrow\quad
		\Phi_1 = \frac{4\pi N I / c}{R_{\text{м}1} + R_{\text{м}2}/2}.
	\]
	Це і є розв'язок. Зверніть увагу: структура виразу така сама, як для електричного кола з послідовно-паралельним з'єднанням опорів. Магнітний опір бічних гілок, з'єднаних паралельно, дорівнює $R_{\text{м}2}/2$, і він додається до опору центрального стрижня послідовно.
\end{Example}

%% --------------------------------------------------------
\section{Магнітний опір повітряного зазору}
%% --------------------------------------------------------

У реальних магнітних системах майже завжди є повітряні зазори --- наприклад, між осердям і якорем електромагніту, у двигунах, у дроселях з підмагнічуванням. Повітряний зазор різко збільшує магнітний опір кола, оскільки проникність повітря $\mu \approx 1$, тоді як у феромагнетика $\mu$ може сягати $10^3$--$10^5$.

Магнітний опір зазору довжиною $\delta$ і площею $S$:
\begin{equation}\label{eq:air_gap_reluctance}
	R_{\text{м,з}} = \frac{\delta}{S}.
\end{equation}
Порівняємо його з опором феромагнітної ділянки тієї ж довжини:
\begin{equation}\label{eq:reluctance_ratio}
	\frac{R_{\text{м,з}}}{R_{\text{м,ф}}} = \frac{\delta/S}{l/(\mu S)} = \mu \frac{\delta}{l}.
\end{equation}
Навіть якщо зазор у $1000$ разів коротший за феромагнітну ділянку, при $\mu = 10^4$ його опір у $10$ разів більший. Саме тому в магнітних колах зазори відіграють вирішальну роль: вони визначають, який струм потрібен для створення заданого потоку.

\begin{Attention}
	Цей ефект добре відомий кожному, хто хоч раз розбирав електромагнітне реле або дверний замок. Коли якір притягнутий до осердя, зазор малий, магнітний опір малий, і для утримання якоря достатньо невеликого струму. Коли ж якір відпущений, зазор великий, і щоб створити той самий потік, потрібен значно більший струм. Тому в момент спрацьовування реле споживає найбільший струм, а в утриманому стані --- значно менший. Це саме та причина, чому в потужних електромагнітах застосовують форсоване збудження: короткочасно подають підвищену напругу, щоб подолати великий опір зазору, а потім знижують її до робочої.
\end{Attention}

%% --------------------------------------------------------
\section{Нелінійність магнітного кола та крива намагнічування}
%% --------------------------------------------------------

Досі ми вважали проникність $\mu$ сталою. Для реальних феромагнетиків це не так: залежність $B(H)$ нелінійна, і магнітний опір $R_{\text{м}} = l/(\mu(H) S)$ залежить від режиму. На рис.~\ref{tikz:magnetization_curve} показано типову криву намагнічування.

\begin{figure}[h!]
	\centering
	\begin{tikzpicture}
		\begin{axis}[
				axis lines=center,
				xlabel={$H$},
				ylabel={$B$},
				x label style={anchor=north west, at={(axis description cs:0.98,0.5)}},
				y label style={anchor=south east, at={(axis description cs:0,0.98)}},
				xmin=0, xmax=10,
				ymin=0, ymax=1.5,
				ticks=none,
				width=10cm,
				height=6cm,
				samples=200
			]
			\addplot[red, thick, domain=0:10] {1.2*(1 - exp(-0.8*x))};
			\node at (axis cs:2.5, 0.45) {лінійна ділянка};
			\node at (axis cs:7.5, 1.25) {насичення};
			\draw[dashed] (axis cs:0, 1.2) -- (axis cs:10, 1.2);
		\end{axis}
	\end{tikzpicture}
	\caption{Типова крива намагнічування феромагнетика. На початковій ділянці $B \approx \mu H$; при великих $H$ настає насичення, і $\mu$ різко падає.}
	\label{tikz:magnetization_curve}
\end{figure}

Через нелінійність закон Ома для магнітного кола \eqref{eq:Ohm_magnetic} не можна застосовувати безпосередньо до всієї кривої. Є два основних підходи:

\begin{enumerate}
	\item \emph{Лінеаризація в робочій точці.} Якщо коло працює в околі деякого потоку $\Phi_0$, можна ввести диференціальний магнітний опір $R_{\text{м,диф}} = d\Phi/d\mathcal{F}$ і розв'язувати задачу для малих відхилень. Цей метод використовують, наприклад, при аналізі підсилювачів на магнітних підсилювачах або при розрахунку індуктивності котушки з підмагнічуванням.
	\item \emph{Графічний метод.} Для кожного значення НС будують криву $\Phi(\mathcal{F})$ шляхом графічного додавання спадів НС на ділянках. Це трудомістко, але дає повну картину. У сучасній практиці такі задачі розв'язують чисельно --- методом скінченних елементів у пакетах типу FEMM або COMSOL.
\end{enumerate}

\begin{Attention}
	Нелінійність магнітного кола --- це не просто математична складність, а фізична реальність, яка має важливі практичні наслідки. По-перше, вона обмежує максимальний потік, який можна «загнати» в осердя: при насиченні подальше збільшення струму майже не збільшує потік. По-друге, вона призводить до спотворення форми струму: якщо до котушки з феромагнітним осердям прикласти синусоїдальну напругу, струм виявиться несинусоїдальним через нелінійність кривої намагнічування. Саме тому в трансформаторах і дроселях виникають вищі гармоніки, а іноді --- явище \emphx{ферорезонансу}, коли нелінійна індуктивність входить у резонанс із ємністю кола. Це ще одна причина, чому розрахунок магнітних кіл вимагає обережності.
\end{Attention}

%% --------------------------------------------------------
\section{Застосування: дроселі, електромагніти, магнітні підсилювачі}
%% --------------------------------------------------------

Поняття магнітного кола --- не абстракція. Воно лежить в основі роботи багатьох пристроїв, які ми вже згадували в попередніх розділах. Розглянемо кілька прикладів.

\paragraph*{Дросель.} Котушка індуктивності з феромагнітним осердям --- це, по суті, магнітне коло, призначене для створення великої індуктивності при малому струмі. Наявність осердя збільшує магнітний потік у $\mu$ разів порівняно з повітряною котушкою, тобто збільшує індуктивність. У СГС зв'язок між індуктивністю та магнітним опором має вигляд
\begin{equation}\label{eq:inductance_reluctance}
	L = \frac{4\pi N^2}{c^2\,R_{\text{м}}},
\end{equation}
що є прямим наслідком закону Ома для магнітного кола разом із означенням $\Phi = \frac{1}{c}LI$. Для дроселя з зазором $R_{\text{м}} = R_{\text{м,ф}} + R_{\text{м,з}}$, і саме зазор визначає, при якому струмі осердя насититься.

\paragraph*{Електромагніт.} Підйомний кран, дверний замок, реле --- усі вони використовують силу притягання феромагнітного якоря до осердя котушки зі струмом. Ця сила виникає через енергію магнітного поля в зазорі: у СГС
\begin{equation}\label{eq:electromagnet_force}
	F = \frac{B^2 S}{8\pi}
\end{equation}
на один зазор (див.~виведення в гл.~\ref{MagEnergy}, приклад із підйомною силою електромагніта). Щоб отримати велику силу, потрібна велика індукція $B$, а отже --- великий потік $\Phi = BS$. Розрахунок котушки електромагніту --- це типова задача про магнітне коло: задаємо потік, знаходимо необхідну НС, потім --- струм і кількість витків.

\paragraph*{Магнітний підсилювач.} Це пристрій, у якому малий струм в одній обмотці керує великим струмом в іншій за рахунок зміни магнітного опору осердя. Якщо осердя підмагнічене близько до насичення, то невелика зміна керуючого струму різко змінює диференціальну проникність, а отже --- індуктивність робочої обмотки. Магнітні підсилювачі використовувалися в системах автоматики до появи напівпровідникових приладів, а в деяких спеціальних застосуваннях --- і досі.

\paragraph*{Трансформатор.} Ми вже детально розглянули трансформатор у гл.~\ref{MagEnergy} як систему зв'язаних електричних контурів. Але на нього можна подивитися й з боку магнітного кола: осердя трансформатора --- це магнітне коло, по якому циркулює потік $\Phi$, створений сумарною НС первинної та вторинної обмоток. У режимі холостого ходу струм первинної обмотки визначається магнітним опором осердя; при навантаженні вторинна обмотка створює свою НС, яка компенсує частину первинної, і потік майже не змінюється. Саме тому струм холостого ходу малий, а струм навантаження може бути великим.

\begin{Attention}
	Погляд на трансформатор як на магнітне коло дозволяє зрозуміти, чому не можна підключати його до джерела постійної напруги. При постійному струмі потік у осерді не змінюється, ЕРС самоіндукції в первинній обмотці дорівнює нулю, і струм обмежується лише активним опором обмотки. Оскільки цей опір малий, струм виявляється величезним --- трансформатор згорає. У колі змінного струму ЕРС самоіндукції врівноважує напругу джерела, і струм залишається в межах норми. Це ще один прояв глибокого зв'язку між електричними та магнітними явищами, який ми простежуємо протягом усієї глави.
\end{Attention}

%% --------------------------------------------------------
\section{Підсумок}
%% --------------------------------------------------------

Ми ввели поняття магнітного кола й показали, що його можна розраховувати тими самими методами, що й електричне коло:

\begin{itemize}
	\item \emph{Закон Ома для магнітного кола:} $\Phi = \dfrac{4\pi\mathcal{F}}{c\,R_{\text{м}}}$, де $R_{\text{м}} = \dfrac{l}{\mu S}$.
	\item \emph{Правила Кірхгофа:} $\sum \Phi_k = 0$ для вузлів; $\sum \Phi_k R_{\text{м}k} = \dfrac{4\pi}{c}\sum \mathcal{F}_i$ для контурів.
	\item \emph{Послідовне й паралельне з'єднання:} магнітні опори додаються при послідовному з'єднанні та складаються обернено при паралельному --- точно як електричні.
	\item \emph{Повітряний зазор:} різко збільшує магнітний опір, і саме він часто визначає режим роботи пристрою.
	\item \emph{Нелінійність:} реальні феромагнетики мають залежну від поля проникність, тому точний розрахунок вимагає або лінеаризації, або чисельних методів.
\end{itemize}

Ці поняття застосовуються в дроселях, електромагнітах, трансформаторах, електричних машинах і магнітних підсилювачах. Вони дозволяють інженеру розраховувати магнітні системи, не розв'язуючи щоразу рівняння Максвелла, --- так само, як закони Кірхгофа дозволяють розраховувати електричні кола в квазістаціонарному наближенні.

Наприкінці варто нагадати головну ідею: магнітне коло --- це не просто зручна аналогія. Це самостійна фізична модель, яка відображає фундаментальну єдність електромагнітних явищ. Струм і потік, ЕРС і НС, опір і релюктанс --- це різні прояви одних і тих самих законів природи, записаних у різних змінних. Розуміння цієї єдності --- ключ до вільного володіння електродинамікою.

%% --------------------------------------------------------
\begin{Questions}
\begin{quizlist}
	\item Що таке магнітне коло? Чим воно відрізняється від електричного кола і що між ними спільного?
	\item Виведіть закон Ома для магнітного кола $\Phi = 4\pi\mathcal{F}/(cR_{\text{м}})$ із теореми про циркуляцію магнітного поля. Яку роль відіграє намагнічувальна сила?
	\item Що таке магнітний опір? Запишіть його через довжину, площу перерізу та проникність. Порівняйте з електричним опором.
	\item Сформулюйте перше й друге правила Кірхгофа для магнітного кола. Як вони пов'язані з аналогічними правилами для електричних кіл?
	\item Розв'яжіть задачу про розгалужене магнітне коло (рис.~\ref{tikz:magnetic_circuit}). Як знайти потоки в гілках, якщо відомі геометрія та НС?
	\item Чому повітряний зазор різко збільшує магнітний опір? Виведіть відношення опорів зазору й феромагнітної ділянки. Як це впливає на роботу електромагніту?
	\item Чому магнітний опір феромагнетика нелінійний? Як виглядає крива намагнічування (рис.~\ref{tikz:magnetization_curve}) і що таке насичення?
	\item Які два підходи використовують для розрахунку нелінійних магнітних кіл? У чому полягає лінеаризація в робочій точці?
	\item Як поняття магнітного кола застосовується до дроселя? Виведіть формулу $L = 4\pi N^2/(c^2 R_{\text{м}})$.
	\item Як розраховується електромагніт? Чому сила притягання залежить від індукції в зазорі?
	\item Що таке магнітний підсилювач? Як він використовує нелінійність кривої намагнічування?
	\item Подивіться на трансформатор як на магнітне коло. Чому струм холостого ходу малий, а струм навантаження може бути великим?
	\item Чому трансформатор не можна підключати до джерела постійної напруги? Поясніть це і з боку електричного кола, і з боку магнітного.
	\item Які практичні наслідки має нелінійність магнітного кола? Що таке ферорезонанс?
	\item Які пристрої використовують магнітні кола? Наведіть приклади з повсякденного життя.
\end{quizlist}
\end{Questions}

\end{document}